% Requires:
% \usetikzlibrary{3d}
% \usetikzlibrary{arrows.meta}
% \usetikzlibrary{calc}
% \usetikzlibrary{decorations.pathmorphing}
% \usetikzlibrary{patterns}
\begin{tikzpicture}[>=latex, line join=round, line cap=round]

  % ==========================================
  % Subfigure (a): Convex 2D cross-section
  % ==========================================
  \begin{scope}

    % Shaded Region

    \begin{scope}
      \pgfmathsetseed{3} % For reproducible randomness
      \fill[pattern=north east lines]
      (0,0)
      % o-face (convex)
      .. controls (-1.5, 0.9) and (-2.5, 1.3) .. (-3, 1.2)
      % Back edge (now prominently convex, bulging outwards to x=-3.4)
      decorate[decoration={random steps,segment length=2mm,amplitude=1mm}] { -- (-3, -1.2) }
      % $n$-face (convex)
      .. controls (-2.5, -1.3) and (-1.5, -0.9) .. (0,0);
      \fill[white] (180+31:0.8) -- (180+31:1.8)
      decorate[decoration={random steps,segment length=2mm,amplitude=1mm}] { -- (180-31:1.8) }
      -- (180-31:0.8)
      decorate[decoration={random steps,segment length=2mm,amplitude=1mm}] { -- cycle};
    \end{scope}

    \pgfmathsetseed{3} % For reproducible randomness
    \draw[thick]
    (0,0)
    % o-face (convex)
    .. controls (-1.5, 0.9) and (-2.5, 1.3) .. (-3, 1.2) coordinate[pos=.1] (start_a0) node[pos=.75, sloped, above] {$0$-face}
    % Back edge (now prominently convex, bulging outwards to x=-3.4)
    decorate[decoration={random steps,segment length=2mm,amplitude=1mm}] { -- (-3, -1.2) }
    % $n$-face (convex)
    .. controls (-2.5, -1.3) and (-1.5, -0.9) .. (0,0) coordinate[pos=.6] (start_an) node[pos=.25, sloped, below] {$n$-face};

    \draw[<-,thick] (start_an) -- ++(-.2,2.5) node[above] {$a_n$};
    \draw[<-,thick] (start_a0) -- ++(-1.2,-2.5) node[below] {$a_0$};

    % Axes
    \draw[->, thick] (0,0) -- (1.0, 0) node[right] {$x$};
    \draw[->, thick] (0,0) -- (0, 1.0) node[above] {$y$};

    % Tangents and Angle
    \draw[dashed, thick, gray, opacity=.5] (0,0) -- (-1.5, -0.9);
    \draw[dashed, thick, gray, opacity=.5] (0,0) -- (-1.5, 0.9);

    \draw[<->, thick] (180+31:1.0) arc (180+31:180-31:1.0);
    \node at (-1.3, 0) {$\alpha$};

  \end{scope}

  % ==========================================
  % Subfigure (b): 3D surface using convex 2D logic
  % ==========================================
  \begin{scope}[
      xshift=4.3cm, yshift=-1cm,
      % Custom 3D perspective to match the view angle
      x={(0.866cm, -0.5cm)}, y={(0cm, 1cm)}, z={(0.866cm, 0.5cm)}
    ]

    % 3D Surface Backings
    \fill[white]
    (0,0,0)
    % Front o-face (z=0)
    .. controls (-1.5, 0.9, 0) and (-2.5, 1.3, 0) .. (-3, 1.2, 0)
    % Top outer sweep curve (z=0 to z=4), adjusted to match the convex crack front
    .. controls (-2.2, 0.8, 1.33) and (-1.4, 0.6, 2.66) .. (-1.5, 0.7, 4)
    % Back o-face
    .. controls (-1.0, 0.8, 4) and (0, 0.4, 4) .. (1.5, -0.5, 4)
    % Edge sweep curve (crack front) back to origin (reversed path, explicitly convex)
    .. controls (1.5, -0.6, 2.66) and (1.0, -0.4, 1.33) .. (0,0,0) coordinate[pos=.5] (mid_front);

    % Front 2D Cross Section at z=0
    \begin{scope}[canvas is xy plane at z=0]

      \pgfmathsetseed{3} % For reproducible randomness
      \fill[pattern=north east lines,draw,thick]
      (0,0)
      % o-face (convex)
      .. controls (-1.5, 0.9) and (-2.5, 1.3) .. (-3, 1.2) node[pos=.3, above right,rotate=-50] {$0$-face}
      % Back edge (now prominently convex, bulging outwards to x=-3.4)
      decorate[decoration={random steps,segment length=2mm,amplitude=1mm}] { -- (-3, -1.2) }
      % $n$-face (convex)
      .. controls (-2.5, -1.3) and (-1.5, -0.9) .. (0,0) node[pos=.25, below,rotate=-20] {$n$-face};
    \end{scope}

    % 3D Surface Boundaries
    % Top outer edge sweep (following convex shift)
    \pgfmathsetseed{3} % For reproducible randomness
    \draw[thick] (-3, 1.2, 0) decorate[decoration={random steps,segment length=2mm,amplitude=1mm}] { -- (-1.5, 0.7, 4) coordinate[pos=.5] (mid_back)};

    \draw[<-,thick] (mid_front) -- ($(mid_front)!.5!(mid_back)$) node[midway, above,sloped] {$a_e$};

    % Far back o-face curve
    \draw[thick] (1.5, -0.5, 4) .. controls (0, 0.4, 4) and (-1.0, 0.8, 4) .. (-1.5, 0.7, 4);

    % The Crack Front (Diffraction Edge) - now prominently convex
    %\draw[thick] (0,0,0) -- (1.5, -0.5, 4);
    \draw[thick] (0,0,0) .. controls (1.0, -0.4, 1.33) and (1.5, -0.6, 2.66) .. (1.5, -0.5, 4);
  \end{scope}

\end{tikzpicture}
